% ID: SEL-D0075-OA-E03
% ESTADO: cotejada. Segunda pasada visual y borrador compilado revisados.
% Fuente: Selectividad/sel_2008_matematicas_aplicadas/Mat CCSS 3 - Examen y criterios de correción.pdf
% SHA-256: ebf609a5063d63be15b900270973808e3fe0a239a4d911d6b41d80544711aeae
% Página: 1; opción A; ejercicio original 3.

\noindent\underline{Parte I}
\begin{enumerate}[label=\alph*)]
\item \textbf{(1 punto)} Sean $A$ y $B$ dos sucesos de un mismo espacio muestral. Sabiendo que $P(A)=0.5$, que $P(B)=0.4$ y que $P(A\cup B)=0.8$, determine $P(A/B)$.
\item \textbf{(1 punto)} Sean $C$ y $D$ dos sucesos de un mismo espacio muestral. Sabiendo que $P(C)=0.3$, que $P(D)=0.8$ y que $C$ y $D$ son independientes, determine $P(C\cup D)$.
\end{enumerate}
\par\noindent\begin{minipage}{\linewidth}
\noindent\underline{Parte II}

El número de días de permanencia de los enfermos en un hospital sigue una ley Normal de media $\mu$ días y desviación típica 3 días.
\begin{enumerate}[label=\alph*)]
\item \textbf{(1 punto)} Determine un intervalo de confianza para estimar $\mu$, a un nivel del $97\%$, con una muestra aleatoria de 100 enfermos cuya media es $8.1$ días.
\item \textbf{(1 punto)} \textquestiondown{}Qué tamaño mínimo debe tener una muestra aleatoria para poder estimar $\mu$ con un error máximo de 1 día y un nivel de confianza del $92\%$?
\end{enumerate}
\end{minipage}\par
